\documentclass[12pt]{article} % Clase de documento: artículo y tamaño de letra
\usepackage[utf8]{inputenc} % Escritura en castellano con acentos
\usepackage[T1]{fontenc} % Escritura en castellano con acentos
\usepackage{calligra} 
\usepackage{pslatex}  % Fuente de letras
\usepackage{graphicx}
\usepackage[margin=2.5cm]{geometry}
\usepackage{txfonts}
\usepackage{xcolor}
\usepackage{fancyhdr}
\usepackage{graphicx}
\usepackage{float}
\usepackage{listings}
\usepackage{multicol}
\usepackage{color}
\usepackage{times}
\usepackage{tikz}
\usepackage{verbatim}
\usepackage{mdwlist}
\usepackage{array}
\usepackage{caption}

\usetikzlibrary{chains,fit,shapes,arrows,calc,shapes,decorations.pathreplacing}


\pagestyle{fancy}
\headheight=50pt %para cambiar el tamaño del encabezado
\fancyhead[L]    %la "L" indica a la izquierda
{	
 \begin{minipage}{3.3cm}
  \includegraphics[width=1.0\textwidth]{Logo-UTN-BA.jpeg}
%  \includegraphics[width=0.6\textwidth]{UTNlogo.eps}
 \end{minipage}	
 \begin{minipage}{7.7cm}
   \normalsize
   {
     \textsf
     {
       \calligra{Universidad Tecnológica Nacional\\ Facultad Regional Buenos Aires \\ Departamento de {Ingeniería} {Electrónica}} 
     }
   }
 \end{minipage}
}

\fancyhead[R] %la "R" indica a la derecha
{
  \begin{minipage}{4.0cm}
   \small
   {
     \emph{\textbf{Informática I}} \\ \emph{\today} \\ \emph{Ejemplos de parcial}\\ \emph{Curso r1092}  
   }
  \end{minipage}}


\begin{document} % Inicio del documento
\newpage
%Cambia de página, el texto después de este comando aparecerá en la siguiente página en adelante.


\noindent
  \begin{center}
   \begin{tabular}{| c | c | c |}
    \hline
     Apellido y Nombres \hspace{8cm} &  Legajo & {Calificación} \\ \hline 
      &	& \\ \hline
   \end{tabular}	
  \end{center}

\noindent
Todas las funciones pedidas deben ser colocadas en el archivo {\bf primerParcialFunciones.c} y debe estar acompañado por su correspondiente {\bf punto h}.\\
Además se debe escribir un programa de testeo para cada función requerida. El nombre de esos archivos es {\bf primerParcialNumeroEjercicio.c}.\\
\\
%%%%%%%%%%%%%%%%%%%%%%%%%%%%%%%%%%%%%%%%%%%%%%%%%%%%%%%%%%%%%%%%%%%%%%%%%%%%%%%%%%%%%%%%%%%%%%%%%%%%%%%%%%%%% Comienzo del tema del examen
%%%%%%%%%%%%%%%%%%%%%%%%%%%%%%%%%%%%%%%%%%%%%%%%%%%%%%%%%%%%%%%%%%%%%%%%%%%%%%%%%%%%%%%%%%%%%%%%%%%%%%%%%%%%
\lstset{
	frame=Ltb,
	framerule=0pt,
	aboveskip=0.5cm,
	framextopmargin=3pt,
	framexbottommargin=3pt,
	framexleftmargin=0.4cm,
	framesep=0pt,
	rulesep=.4pt,
	backgroundcolor=\color{gray97},
	rulesepcolor=\color{black},
% 	language=C,
	captionpos=b,
	tabsize=3,
	frame=lines,
	keywordstyle=\color{blue},
	commentstyle=\color{Gray},
	stringstyle=\color{red},
	numbers=left,
	numberstyle=\tiny,
	numbersep=5pt,
	breaklines=true,
	showstringspaces=false,
	basicstyle=\ttfamily\scriptsize,
	emph={label},
	framerule=0pt,
}

Implemente una función que realice la validación de una CBU (Clave Bancaria Uniforme).\\
La CBU está formado de la siguiente manera.\\
\fontsize{9pt}{9pt}\selectfont
\vspace{-0.6cm}
\begin{center} 
   \begin{tabular}{|c|c|c|c|c|c|c|c|c|c|c|c|c|c|c|c|c|c|c|c|c|c|}
    \hline 
     0&1&4&0&1&2&5&6&5&5&6&5&4&1&8&5&4&7&6&5&4&3 \\ \hline 
     $E_0$&$E_1$&$E_2$&$S_0$&$S_1$&$S_2$&$S_3$&$D_0$&$C_0$&$C_1$&$C_2$&$C_3$&$C_4$&$C_5$&$C_6$&
     $C_7$&$C_8$&$C_9$&$C_{10}$&$C_{11}$&$C_{12}$&$D_1$\\ \hline 
   \end{tabular}
\end{center}   	
%\vspace{2.5cm}
\fontsize{12pt}{12pt}\selectfont
Donde:
\begin{itemize}
\item {\bf E0 a E2}: Es el número de la entidad bancaria.
\item {\bf S0 a S3}: Es el número de sucursal.
\item {\bf D0}: Es el dígito verificador de E y S
\item {\bf C0 a C12}: Es el número de cuenta.
\item {\bf D0}: Es el dígito verificador de C
\end{itemize}
Los dígitos verificadores de la clave bancaria única se calculan de la siguiente forma:\\
\fontsize{10.4pt}{10.4pt}\selectfont

$R_0$ = $E_0$*9 + $E_1$*7 + $E_2$*1 + $S_0$*9 + $S_1$*7 + $S_1$*1 + $S_2$*3 \\
$D_0$ = 10 - ($R_0$ \% 10)\\

$R_1$ = $C_0$*9 + $C_1$*7 + $C_2$*1 + $C_3$*3 + $C_4$*9 + $C_5$*7 + $C_6$*1 + $C_7$*3 + $C_8$*9 + $C_9$*7 + $C_{10}$*1 + $C_{11}$*3  + $C_{12}$* 9\\
$D_1$ = 10 - ($R_1$ \% 10)\\
\fontsize{12pt}{12pt}\selectfont
\\
La función tiene el siguiente prototipo:
\begin{center}
{\bf int cbu\_validar (char *dataPtr)}
\end{center}
Donde {\bf dataPtr} es el puntero la clave bancaria uniforme a validar terminada en ‘{\bf $\setminus$0}’\\
\\
Devuelve:
\begin{itemize}
\item {\bf -1} Si la CBU pasada no tiene 22 caracteres.
\item {\bf -2} Si alguno de los caracteres de la CBU no es un número.
\item {\bf -3} Si el dígito verificador D0 no corresponde.
\item {\bf -4} Si el dígito verificador D1 no corresponde.
\end{itemize}
 
\end{document} % Fin del documento.
 

  
